\section{FoldPot SNARK}
The starting point of our framework 
is a combination of SuperNova \cite{SuperNova} and CycleFold \cite{CycleFold},
and our implementation is based on deriving Sonobe \cite{Sonobe}, which
implements Nova \cite{Nova} plus CycleFold \cite{CycleFold}.
SuperNova supports non-uniform circuits, and
cyclefold allows us to take advantage of half pairing groups
like BN254/Grumpkin.

Our contribution in the framework is the introduction of a 
cycle pairing component in producing compressed decider proof.
We observe that non-native arithmetics such as Pairing operations for
BN254 can also be folded, by generalizing CycleFold.
This allows us to use Quazi-linear NIZK (QA-NIZK) \cite{Gonzalez15,LegoSnark}
to encode the decider proof for many non-uniform circuits as pairing
operations in circuit.

\subsection{Folding Scheme: SuperNOVA + CycleFold}
Let $\G_1$ and $\G_3$ be two curve groups (e.g., that of
BN254 \cite{BN254} and Grumpkin \cite{Grumpkin}) s.t. the scalar field
of $\G_1$ is the base field of $\G_3$ and 
the base field of $\G_1$ is the scalar field of $\G_3$.
We also assume bilinear pairing is supported on the curve 
for $\G_1$ (e.g., BN254), i.e.  there exists 
$(\G_1, \G_2, \G_T, e)$ and 
$e: \G_1 \times \G_2 \rightarrow \G_T$
is a non-degenerating bilinear pairing function.
We use additive group notations
from Groth16 \cite{Groth16}. 
We let $\F_r$ and $\F_q$ be the scalar and base prime
field of $\G_1$ (i.e., base and scalar field of $\G_3$).
For example, let $a,b \in \F_r$, then
$[ab]_T = e([a]_1, [b]_2)$ denotes that the pairing
of $[a]_1 \in \G_1$ and $[b]_2 \in \G_2$ results in
$[ab]_T \in \G_T$.

\subsubsection{CyclePair Gadget} We first briefly review
CycleFold, based on which we develop CyclePair.
The key observation of CycleFold is that: 
curve operations in $\G_1$ can be represented {\em natively}
over $\G_3$, because $\G_1$'s point in the affine form,
can be represented as two scalar field elements of $\G_3$.
Let $k = 2\lambda$ where $\lambda$ is the security parameter.
Formally, the CycleFold relation is denoted as
$\cyclefold$ over $\F_q$, i.e., $\myvec{x}$ (i/o) and $\myvec{w}$ (witness)
are vectors of $\F_q$ elements:
\[
	(\myvec{x}; \myvec{w}) \in \cyclefold ~\textrm{iff.}~
	\left\{
	\begin{array} {l l}
		& \myvec{x} = ([a]_1, [b]_1, [c]_1, r)  \\
		& [c]_1 = [a]_1 + r [b]_1 \\
		& r \in \{0,1\}^k \\
	\end{array}
	\right.
\]

In folding schemes like Nova \cite{Nova},
the committed instance of $\cyclefold$ can be represented
as an R1CS over $\F_r$. Thus in CycleFold, the system
can fold the committed instances for both the main circuit and
the cyclefold circuit {\em natively} over $\G_3$ and $\G_1$ respectively
at reduced cost.  In the following, we write
$([a]_1, [b]_1, [c]_1, r) \in \cyclefold$ for short (saving the
notation of witness). Similarly, we can use $\cyclepair$
to denote the {\em cycle pair} relation over $\F_q$:
\[
	(\myvec{x};\myvec{w}) \in \cyclepair~\textrm{iff.}~
	\left\{
	\begin{array} {l l}
		& \myvec{x} = ([c]_T, [a]_1, [b]_2, [d]_T)  \\
		& [d]_T = [c]_T + e([a]_1, [b]_2) \\
	\end{array}
	\right.
\]

Here we have one point in $\G_1$ and $\G_2$ respectively and
two points in $\G_T$, all can be represented as elements in 
$\F_q$ and the relation encodes one pairing and one $\G_T$
operation (natively). In practice it costs $12k$ R1CS constraints
over $\F_q$.

\subsubsection{Mixed R1CS} The use of cyclepair gadget allows
encoding a mixed R1CS relation over two fields. We first
give a motivating example for the use of such mixed R1CS. 

In the decider proof for SuperNova, each non-uniform circuit has
a commitment $\overline{\myvec{W}^{(i)}}$ to its witness 
$\myvec{W}^{(i)}$. The prover needs to demonstrate the knowledge
of them, and each circuit has a separate Pedersen vector commitment 
key.\footnote{We note
that the authors of SuperNova mentions in a private communication that
when circuits have the same size, they do can share the same Pedersen
vector commitment keys. Here, when the circuit size are different, 
using separate Pedersen commitment keys for individual circuits
avoids the extra proof to justify the degree bound of
commitments.} It is possible to batch-prove all $\myvec{W}^{(i)}$
using QA-NIZK. However, we would like to perform the QA-NIZK verification
in-circuit to further compress proof size, given that there is a large
number of circuits.  This would require
{\em encoding bilinear pairing operations in circuit}, if performed
non-natively over $\G_1$, will be very costly.  We show how to
accomplish it with mixed R1CS relation and $\cyclepair$.

We know that any field element $x \in \F_q$ can be represented as 
a vectof of field elements in $\F_r$.\footnote{In practice, to support
the R1CS representation non-natively, we set the limb size to $55$,
inheriting the setting in Sonobe \cite{Sonobe}, so that the lazy evaluation
of non-native representation can be applied \cite{JSnark}.}
We denote it as $\frproj{x}$, and naturally apply it to a vector
of $\F_r$ elements, and similarly to $\G_1$, $\G_2$, and $\G_T$
elements (as they are can be represented as vectors of $\F_q$ elements).

\begin{definition} A {\em mixed R1CS relation} over $\F_r$ and $\F_q$
is a tuple $(\myvec{x}^{(1)} \in \F_r^{m_1}, \myvec{x}^{(2)} \in \F_r^{m_2}, 
\myvec{w}^{(1)} \in \F_r^{n_1}, \myvec{w}^{(2)} \in \F_q^{n_2}, R_1, R_2)$
for some $m_1, m_2, n_1, n_2 \in \Nat$ where
$R_1$ and $R_2$ are two R1CS relations over $\F_r$ and $\F_q$
respectively, and $(\myvec{x}^{(1)} \concat \frproj{\myvec{x}^{(2)}}; \myvec{w}^{(1)}) \in R_1$ and $(\myvec{x}^{(2)}; \myvec{w}^{(2)}) \in R_2$.
\end{definition}

In this paper, we restrict $R_2$ to be $\cyclepair$, but it can be generalized
to any R1CS relation over $\F_q$.  In the following, we show 
$\circqa$, an instance of mixed R1CS relation for capturing
the verification of QA-NIZK proofs.

\begin{definition} $\circqa$ is a mixed R1CS relation:
\[
\left( 
  \begin{array}{l}
  	\left(h^{(a)}, h^{(b)}, h^{(a')}, h^{(b')}\right), \\ 
    \left([c]_T, [a]_1, [b]_2], [d]_T\right), \\
	R_1, \cyclepair
  \end{array}
\right)
\]
such that: $([c]_T, [a]_1, [b]_2, [d]_T) \in \cyclepair$ and
$R_1$ defines the following relation over $\F_r$:
\[
	h^{(a')} = \hash(h^{(a)}, \frproj{[a]_1})  \And
	h^{(b')} = \hash(h^{(b)}, \frproj{[b]_1}) 
\]
\end{definition}

One can chain many instances $\circqa$ as step
circuits to verify QA-NIZK relation below: 
\begin{eqnarray}
	& &\sum_{i=1}^n e([\myvec{x}_i]_1, [\myvec{C}_i]_2) = e(\prfnizk, [a]_2) \\
	&\And& h^{(x)} =\hashchain(\myvec{x}) \\
	&\And& h^{(C)} =\hashchain(\myvec{C})
\end{eqnarray}

Here:
QA-NIZK proof $\prfnizk$ is a $\G_1$ element,
its verification key $\left([\myvec{C}_i]_2, [a]_2\right)$ is a vector of $\G_2$
elements. The I/O of the folded relation are the two hashchains
($h^{(x)}$ and $h^{(C)}$), and they can be further hided in circuit.
When applying to the decider proof of SuperNova, the
$\myvec{x}$ are the $\G_1$ points contained in each committed R1CS instance.

\subsubsection{SuperNova \cite{SuperNova} + CycleFold \cite{CycleFold} with CyclePair}
We now present a slightly refactored SuperNova \cite{SuperNova} + 
CycleFold \cite{CycleFold}
framework to support cyclepair, which also incorporates
the support for finger-printing technique of fixed memory segement
introduced in Section \ref{cf_zksnark}. We use the
$\nifsp$ and $\nifsv$ functions from Nova \cite{Nova} directly,
which folds the concrete and committed R1CS instances.
We first introduce a slightly
modified definition of committed and relaxed R1CS instance.

% As Nova \cite{Nova}, SuperNova \cite{SuperNova}, we define R1CS, extended
% R1CS, and the committed instance below, by slightly adjusting
% instnce I/O and witness order in the original deifnition.
% 
% \begin{definition} \cite[Definitions 8 and 9] {SuperNova} Given finite field 
% $\F$, let $R = (n, m, \myvec{A}, \myvec{B}, \myvec{C})$ 
% denote the structure of an R1CS relation, where
% $\myvec{A}, \myvec{B}, \myvec{C} \in \F^{n\times n}$, each with $\Omega(n)$ non-zero entries,
% and $m$ is the public input length. 
% We say $(\myvec{x};\myvec{w}) \in R$ if and only if
% $(\myvec{A} \cdot \myvec{Z}) \circ (\myvec{B} \cdot \myvec{Z}) = 
% \myvec{C} \cdot \myvec{Z}$ where
% $\myvec{x}\in \F^m$ and $\myvec{W} \in \F^{n-m-1}$ and $\myvec{Z} = 
% (\myvec{W},\myvec{x},1)$.
% 
% Let $\pk_E$ and $\pk_W$ be Pedersen vector keys of vector length
% $n$ and $n-m-1$ respectively. A {\em commited instance} of
% $R$ is a tuple $(\overline{\myvec{E}}, u, \overline{\myvec{W}}, \myvec{x})$
% where $\overline{\myvec{E}} = \com(\myvec{E}; r_E)$ and
% $\overline{\myvec{W}} = \com(\myvec{W}; r_W)$ for
% $(\myvec{E}, r_E, \myvec{W}, r_W) \in (\F^{n}, \F, \F^{n-m-1}, \F)$
% s.t. $(\myvec{A}\cdot \myvec{Z})\circ (\myvec{B} \cdot \myvec{Z}) = u(\myvec{C}\cdot \myvec{Z}) + \myvec{E}$ where $\myvec{E} = (\myvec{W}, \myvec{x}, u)$.
% \end{definition}

% The cyclefold and cyclepair use the standard definition of
% R1CS over $\F_q$ and its committed instance. For the main relation
% over $\F_r$, to apply finger printing techniques, in Definition 
% \ref{def:extr1cs}, an additional fixed witness section is introduced.
% We repeat it here. The difference is a fixed mapping funciton $\rho$
% that defines a fixed region $\myvec{W^{\rho}}$ in $\myvec{W}$.

\begin{definition} (extending \cite[Definition 12]{Nova})
Given finite field $\F$, the structure of a 
{\em Fixed-memory R1CS} (F-R1CS) relation $R$ is denoted by 
$R = (n, m, t, \myvec{A}, \myvec{B}, \myvec{C}, \rho)$ 
where $\myvec{A}, \myvec{B}, \myvec{C} \in \F^{n\times n}$, 
each with $\Omega(n)$ non-zero entries,
and $m$ is the public input length.  $\rho$ is 1-to-1 map (projection) of 
size $t$ s.t. $\forall i\in [t]: \rho(i) \in [0,n-1]$.
Given $\myvec{W} \in \F^n$, 
the projected word $\rho(\myvec{W})$, also 
denoted as $\myvec{W^{(\rho)}}$ has length $t$ and
satisfies: $\forall i \in [t]: \myvec{W^{(\rho)}}_i = \myvec{W}_{\rho(i)}$.

We say $(\myvec{x};\myvec{w}) \in R$ if and only if
$(\myvec{A} \cdot \myvec{Z}) \circ (\myvec{B} \cdot \myvec{Z}) = 
\myvec{C} \cdot \myvec{Z}$ where
$\myvec{x}\in \F^m$ and $\myvec{W} \in \F^{n-m-1}$ and $\myvec{Z} = 
(\myvec{W},\myvec{x},1)$. 
A committed R1CS instance for $R_\rho$ is a tuple 
$(\overline{\myvec{E}}, u,
\overline{\myvec{W}}, \myvec{x}, \overline{\myvec{\myvec{W}^{(\rho)}}})$
where $\overline{\myvec{W}^{(\rho)}}$,
$\overline{\myvec{W}}$, $\overline{\myvec{E}}$ are Pedersen vector
commitment of $\myvec{W^{(\rho)}}$, $\myvec{W}$ and $\myvec{E}$,
s.t. $(\myvec{A}\cdot \myvec{Z})\circ (\myvec{B} \cdot \myvec{Z}) = u(\myvec{C}\cdot \myvec{Z}) + \myvec{E}$ where 
$\myvec{Z} = (\myvec{W}, \myvec{x}, u)$.
\end{definition}

To support the finger-printing techniques, the main circuit (over $\F_r$)
is in F-R1CS, and the cyclefold and cyclepair circuits are the
standard R1CS as in SuperNova. Their committed instsances are
the ones given in SuperNova, i.e., without the $\overline{\myvec{W}^{(\rho)}}$
component.

We present the augmented circuit $\augf{j}$ in  Figure\ \ref{fig_supernova},
with slight refactoring of the algorithm in \cite[Construction 1]{SuperNova}.
Let $\ell$ be the number of circuits.
Each $\augf{j}$ takes a collection of non-deterministic
hints from the prover, and mainly performs two operations: (1) computes
$\z_{i+1} \leftarrow \func{j}(\z_i, \omega_i)$, and (2)
generates $\U_{i+1}$ by folding
$\U_{i}$ (current folded instance) and $\u_{i}$ (incoming instance).
Compared with the standard SuperNova,  
there are three relations that are being folded: the 
main circuit $\func{j}$ over $\F_r$, and the cyclefold and cyclepair
relations over $\F_q$. To support IVC, the circuit 
outputs three $\F_r$ elements ($\xmain$, $\xcyclefold$, $\xcyclepair$), 
which are the hash of the committed instances of these three folded relations. 
Function $\genCycleFoldX()$ generates the input for the cycle fold
instance (the random combination of the $\overline{\myvec{W}}$,
$\overline{\myvec{E}}$, $\overline{\myvec{W}^{(\rho)}}$
for folding). Notice that all outputs are converted to limbs of $\Fr$.

% The workflow of the $\augf{j}$ function simulates that of the SuperNova
% framework \cite[Section 4.2]{SuperNova}, with some minor adjustifications:
% (1) the $\phi(\pc_i, \pc_{i+1})$ simply asserts the validity of
% the choice of $\pc_{i+1}$ (as in our context $\pc_{i+1}$ can be any value
% in $[\ell]$, the prover just selects the least expensive one).
% (2) Folding operations are added for cyclefold and cyclepair relations.
% We re-use most of the notations in \cite{SuperNova}, e.g.,
% $(\ubot{i}, \wbot{i})$ are the trivial satisfying
% instances for circuit $\func{i}$.
% 
% In Figure\ \ref{fig_supernova}, function
% $\nifsv$ performs the folding of committed instances
% for cyclefold and cyclepair. Notice that the commitments
% are in $\G_3$ (where the affine representation are encoded using
% $\F_r$ elements). Function $\rohash()$ represents the
% random oracle function and it maintains an implicit
% transcripts of all past prover/verifier communication.
% 
% Augmented function $\augf{j}$ takes non-deterministic hints
% from the prover (as part of the witness): $\advicemain$, 
% $\advicefold$, and $\advicepair$ for the main relation, and
% the cyclefold and cyclepair relations, respectively.
% We briefly define their semantics below:
% 
% \begin{enumerate}
% 	\item $\advicemain:$ $\pc_i$ is the circuit ID for the $i$'th step.
% 		$\z_i$ is the output of $i$'th step. $\omega_i$ is the
% 		hint (witness) for $\func{j}$.  $\U_i$ is
% 		the folded instances up to step $i$.
% 		$\u_i$ is the incoming instance.
% 		$\overline{T}$ is the computed commitment to error terms 
% 		used in folding committed instances.
% 	\item $\advicefold:$  $\Uc{i}$ is the
% 		folded instance for cyclefold fold,
% 		and $\uc{i}$ the incoming instance,
% 		for step $i$. As for each step, there
% 		are three cyclefold relations (for $\overline{\myvec{E}}$, 
% 			$\overline{\myvec{W}}$, $\overline{\myvec{W}^{(\rho)}}$).
% 		The corresponding commitments to the witness of
% 		these three instances of cyclefold instances are provided,
% 		namely: $\comcw$, $\comce$, and $\comcf$. Similarly
% 		the corresponding committed error terms are
% 		$\comtw$, $\comte$, and $\comtf$.
% 		Notice that folding operation of
% 		$\Uc{i}$ and $\uc{i}$ can be directly computed
% 		over $\F_r$.
% 	\item $\advicepair:$ all hints are similar to that
% 	of $\advicefold$.  The cyclepair input
% 	$([t_1]_T, [a]_1, [b]_2, [t_2]_T)$ is retrieved from
% 	the problem statement of the cyclepair instance, i.e.,
% 	formally each circuit $\func{j}$ is a mixed R1CS relation
% 	and the cyclepair input is retrieved from the input statement
% 	$\myvec{x}^{(2)}$ from its relation $R_2$. 
% 	Note that
% 	it can be serialized to a vector of $\F_r$ elements (using the limb 
% 	approach).
% \end{enumerate}

% The workflow of the $\augf{j}$ in Figure\ \ref{fig_supernova}
% mainly follows that of \cite[Construction 1]{SuperNova},
% by adding the CycleFold \cite{CycleFold}, where the curve point
% operation is out-sourced to three foldings of cyclefold relation.
% Following the same strategy, we adds a folding of CyclePair relation,
% as highlighted. Note that these foldings are performed at $\F_r$,
% and all inputs of relations are projected to $\F_r$ when necessary.
% Finally, it returns the hash of the three relations to support IVC.



\def\hp{\hspace*{0.1in}}
\begin{small}
\begin{algorithm}[bt]
\SetKwData{Left}{left}\SetKwData{This}{this}\SetKwData{Up}{up}
%\SetKwFunction{Union}{Union}\SetKwFunction{FindCompress}{FindCompress}
%$ \U' \leftarrow \nifsv(\vk, \U, \u, \overline{T})$ $\#$ over $\F_r$
%\newline
%(1) Assert $\u.u = 1$ and $\u.\overline{E} = [0]_3$.\newline
%(2) $r \leftarrow \rohash(\vk, \U, \u, \overline{T})$.\newline
%(3) Return $\U' = (
%			\U.\overline{E} + r \overline{T},
%			\U.\overline{\myvec{W}} + r \overline{\myvec{w}},
%			\U.u + r,
%			\U.\myvec{x} + r \myvec{x})$.
%	\\
\SetKwFunction{Union}{Union}\SetKwFunction{FindCompress}{FindCompress}
$ \left(\vecx{1},\vecx{2},\vecx{3}\right) 
	\leftarrow \genCycleFoldX(\vk, \U, \u, \U', \overline{T})$ 
\newline
(1) $r \leftarrow \rohash(\vk, \U, \u, \U', \overline{T})$ \newline
(2) Return $\left(
	\begin{array}{l}
		\frproj{\left(
		\U.\overline{\myvec{W}},
		\u.\overline{\myvec{W}},
		\U'.\overline{\myvec{W}},
		r\right)} \\
		\frproj{\left(
			\U.\overline{\myvec{E}},
			\u.\overline{\myvec{E}},
			\U'.\overline{\myvec{E}},
			r\right)} \\
		\frproj{\left(
		\U.\overline{\myvec{W^{(\rho)}}},
		\u.\overline{\myvec{W^{(\rho)}}},
		\U'.\overline{\myvec{W^{(\rho)}}},
		r\right)}
	\end{array}
	\right)$.


\SetKwFunction{Union}{Union}\SetKwFunction{FindCompress}{FindCompress}
$ \left(\xmain,\xcyclefold,\xcyclepair\right) 
	\leftarrow \augf{j} \left(\vk, \left(
		\advicemain, \advicefold, \advicepair\right) \right)$
\newline
(1) Parse  \newline
	$\advicemain$ $  =
		\left(\U_i, \u_i, \U_{i+1}, \pc_i, \pc_{i+1}, (i, \z_0, \z_i),  
			\omega_i, \O_i, \overline{T}\right)$.  \newline
		$\advicefold$   $= \left(
				\begin{array}{l}
					\Uc{i}, \Uc{i+1}, \uce{1}, \uce{2}, \uce{3}, \\
					\comCT{1}, \comCT{2}, \comCT{3}
				\end{array}
		  \right).
		$   \newline
	$\advicepair$  	
		$= \left(
		   \Upair{i}, \Upair{i+1}, \upair{i}, \comPT, \xinpair
		  \right)$  \newline
(2) Run $\phi(\pc_i, \pc_{i+1})$ to assert validty of $\pc_{i+1}$.\newline
(3) $\z_{i+1} \leftarrow \func{j}(\z_i, \omega_i)$.\newline 
(4) Assert $\z_i = \hash(\O_i)$. \newline
(5) Assert $\O_i.\CpInput = \xinpair$. \newline
(6) If $i=0$: 
	assert $\z_i = \z_0$, and \newline
	$\htab \left(
			\U_{i+1}, 
			\Uc{i+1}, 
			\Upair{i+1}
			\right) \leftarrow 
		\left(\left(\ubot{i}\right)_{i=1}^\ell, \ucbot, \upbot \right)$ 
	\newline
(7) If $i>0$:   \newline
$\htab$ (7.1) Check $\u_i$ is the commitment to a standard R1CS, i.e.,
		  $\u_i.x$ $=$ $\hash(\vk,$ $i, \pc_i,$ $\z_0, \z_i, \U_i)$,
		  and $\u_i.\overline{E} = \ubot{i}.\overline{E}$,
		  and $\u_i.u = 1$. \newline
$\htab$(7.2) Check for $k\neq \pc_i$: $\U_{i+1}[k] = \U_i[k]$ and \newline
$\htab$(7.3) Perform folding of $\cyclefold$ and $\cyclepair$ natively
	over $\F_r$: \newline
	\newline
	$\htab\htab$ $\Ucnop \leftarrow \nifsv\left(\vk, \Uc{i}, \ucnopone, \comtw\right)$ \newline
	$\htab\htab$ $\Ucnop \leftarrow \nifsv\left(\vk, \Ucnop, \ucnoptwo, \comte\right)$ \newline
	$\htab\htab$ $\Uc{i+1} \leftarrow \nifsv\left(\vk, \Ucnop, \ucnopthree,  \comtf\right)$ \newline
	\colorbox{yellow}{$\htab\htab \Upair{i+1} \leftarrow \nifsv\left(\vk, \Upair{i}, \upair{i}, \comtwp\right)$} \newline
$\htab$(7.4) Assert Inputs: \newline
	$\htab\htab \left( \uce{i}.\myvec{x} \right)_{i=1}^3 = 
		\genCycleFoldX\left(
			\begin{array}{l}
				\vk, \U_i[\pc_i], \u_i,  \\
				\U_{i+1}[\pc_i], \overline{T}
			\end{array}
		\right)$ \newline
	\colorbox{yellow}{
		$\htab\htab$ $\upair{i}.\myvec{x} = \xinpair$
	} \newline 
(8) Return 
$\left(
\begin{array}{l}
    \hash\left(\vk, (i, \z_0, \z_{i+1}), \U_{i+1}\right),  \\
	\hash\left(\Uc{i+1}\right),  \\
	\hash\left(\Upair{i+1}\right)
\end{array}
\right)$.

\caption{Modified Augmented Function from SuperNova \cite{SuperNova} 
	and CycleFold \cite{CycleFold}}
\label{fig_supernova}
\end{algorithm}
\end{small}

In Figure\ \ref{fig_folding},
we then provide a refactored IVC framework 
presented \cite[Construction 1]{SuperNova}, by adding cyclepair.
The $\ProveIVC()$ function prepares the
input for $\augf{j}$, and generates the next IVC proof
which consists of the pairs of folded and incoming
R1CS instances.  Here $\nifspwrapper$ is a wrapper of $\nifsp$, 
which given a committed instance pair
$(\U_i, \W_i)$ for a relation $R$, an input $\myvec{x}$ for $R$,
finds the corresponding witness for $\myvec{x}$ and generates
an incoming instance $(\u_i, \w_i)$ and fold it with $(\U_i, \W_i)$.
We know that for $\cyclepair$ and $\cyclefold$ relations,
it takes linear time to compute the witness for a given input.
To minimize the public I/O size, each $\z_i$ is 
generated from a fixed structure $\O_i$ by
$\hash(\O_i)$. In $\O_i$ there is a field of $\CpInput$
of form $\left([t_1]_T, [a]_1, [b]_2, [t_2]_T\right)$ which serves as
the input for $\cyclepair$. For arithmetic function $\func{j}$
we use $\realfunc{j}\left(\pk, \myvec{x}, \omega_i\right)$ to 
output the $\O_{i+1}$ instance, thus
$\func{j}$ can be written as $\hash\left(\realfunc{j}\left(\pk, \O_i, \omega_i\right)\right)$. 


\def\hp{\hspace*{0.1in}}
\begin{small}
\begin{algorithm}[bt]
\SetKwData{Left}{left}\SetKwData{This}{this}\SetKwData{Up}{up}
\SetKwFunction{Union}{Union}\SetKwFunction{FindCompress}{FindCompress}
$ (\U_{i+1}, \W_{i+1}, \u_{i}, \w_{i}, \comT) \leftarrow \nifspwrapper(\pk, \U_i, \W_i, \vecxnop, R)$ \newline
	(1) Given $\vecxnop$ and $R$, compute $(\vecxnop; \vecwnop) \in R$. \newline
	(2) Let $\u_{i}$ be committed R1CS instance of $(\vecxnop; \vecwnop)$
		for $R$. \newline
	(3) $(\U_{i+1}, \W_{i+1}, \comT)  \leftarrow \nifsp\left(\pk, (\U_i, \W_i),
		(\u_i, \w_i)\right)$. \newline
	(4) Return $(\U_{i+1}, \W_{i+1}, \u_i, \w_i, \comT)$.
\newline

\SetKwFunction{Union}{Union}\SetKwFunction{FindCompress}{FindCompress}
$ \left( (\pk_i, \vk_i) \right)_{i=1}^{\ell} \leftarrow \SetupIVC\left(\lambda, \{F_i\}_{i=1}^\ell \right)$. \cite[pp. 17]{SuperNova}\newline
(1) $\pp_i \leftarrow \nifsg(1^\lambda)$ for each $i \in [\ell]$. \newline
(2) Let $(\pk_i, \vk_i)  \leftarrow \nifsk(\pp_i, S(\augf{i}))$ for
	each $i \in [\ell]$. \newline
(3) Output $( (\pk_i, F_i) )_{i=1}^{\ell}$ as $\pk$ and
	$ ( (\vk_i, F_i) )_{i=1}^{\ell}$ as $\vk$.

\SetKwFunction{Union}{Union}\SetKwFunction{FindCompress}{FindCompress}
$ \ivcprf_{i+1} \leftarrow \ProveIVC(\pk, \vk, (i, z_0, z_i), \omega_i, \pc_{i+1}, 
	\ivcprf_i)$. \newline
(1) Parse $\ivcprf_i$ as $\left(
%	\begin{array}{l}
		\left(\U_i, \W_i, \u_i, \w_i\right), 
		\left(\Uc{i}, \Wc{i}\right), 
		\left(\Upair{i}, \Wpair{i}\right),
		\pc_i, \O_i
%	\end{array}
	\right)$. \newline
(2) $\left(\xin{fold}{1},\xin{fold}{2},\xin{fold}{3}\right) 
		\leftarrow \genCycleFoldX(\vk, \U_i, \u_i) $ \newline
(3) $\xinpair \leftarrow \O_i.\CpInput$. \newline
	\newline
(5) $\O_{i+1} \leftarrow \realfunc{\pc_{i+1}}\left(\O_i, \omega_i\right)$ and
	$\z_{i+1} \leftarrow \hash(O_{i+1})$. \newline
(6) Compute:  \newline
	$\left(\big(\Uce{j},\Wce{j}, \uce{j}, \comCT{j}\big)\right)_{j=1}^3 \leftarrow$ 
	\newline
	$\htab \left(\nifspwrapper(\pk, \Uc{i}, \Wc{i}, \xin{fold}{j}, \cyclefold)\right)_{j=1}^3$.
	\newline
	$\left(\Upair{i+1},\Wpair{i+1}, \upair{i}, \comPT\right) \leftarrow$ 
	\newline
	$\htab \nifspwrapper(\pk, \Upair{i}, \Wc{i}, \xinpair, \cyclepair)$.
	\newline
(7) Let  $\advice = \left(\advicemain, \advicefold, \advicepair\right)$ 
	where
	\newline $\advicemain$ $  =
		\left(\U_i, \u_i, \U_{i+1}, \pc_i, \pc_{i+1}, (i, \z_0, \z_i),  
			\omega_i, \O_i, \overline{T}\right)$.  \newline
		$\advicefold$   $= \left(
				\begin{array}{l}
					\Uc{i}, \Uc{i+1}, \uce{1}, \uce{2}, \uce{3}, \\
					\comCT{1}, \comCT{2}, \comCT{3}
				\end{array}
		  \right).
		$   \newline
	$\advicepair$  	
		$= \left(
		   \Upair{i}, \Upair{i+1}, \upair{i}, \comPT, \xinpair
		  \right)$  \newline
(8) $(\u_{i+1}, \w_{i+1}) \leftarrow \trace\left(\augf{\pc_{i+1}}\left(\vk, \advice\right)\right)$ \newline
(9) Output $\ivcprf_{i+1}$ as $\left(
	\begin{array}{l}
		\left(\U_{i+1}, \W_{i+1}, \u_{i+1}, \w_{i+1}\right), \\
		\left(\Uc{i+1}, \Wc{i+1} \right), 
		\left(\Upair{i+1}, \Wpair{i+1}\right),
		\pc_{i+1}
	\end{array}
	\right)$. \newline

\SetKwFunction{Union}{Union}\SetKwFunction{FindCompress}{FindCompress}
$ (0,1) \leftarrow \VerifyIVC\left(\vk, (i, z_0, z_i), \ivcprf_i\right)$. \newline
If $i=0$ check $z_i = z_0$ Otherwise Parse
$\ivcprf_{i}$ as $\left(
	\begin{array}{l}
		\left(\U_{i}, \W_{i}, \u_{i}, \w_{i}\right), 
		\left(\Uc{i}, \Wc{i}\right), 
		\left(\Upair{i}, \Wpair{i}\right),
		\pc_{i} 
	\end{array}
	\right)$. \newline
check the following: \newline
(1) $1 \leq \pc_i \leq \ell$. \newline
(2)  Assert $\u_i.\myvec{x} = 
		\left(
		\begin{array}{l}
			\hash\left(\vk, (i, \z_0, \z_{i}), \U_{i}\right),  \\
			\hash\left(\Uc{i}\right),  \\
			\hash\left(\Upair{i}\right)
		\end{array}
	\right)$. \newline
(3) $\forall i\in [\ell]$:  $\u_{i}$, $\uc{i}$, $\upair{i}$ are standard, i.e.,
\newline
	their $\overline{\myvec{E}}$ is $\ubot.\overline{\myvec{E}}$
	and $u = 1$. \newline
(4) Assert:
$\left\{\begin{array}{l}
	(\u_i, \w_i) \in \augf{\pc_i} \\ 
	(\Uc{i}, \Wc{i}) \in \cyclefold \\
	(\Upair{i}, \Wpair{i}) \in \cyclepair  \\
 \end{array}
\right.$ \newline
(5) Assert: $\forall j\in [\ell]$, $(\U_{i}[j], \W_{i}[j]) \in \augf{j}$.  
\newline

\caption{Main IVC Framework}
\label{fig_folding}
\end{algorithm}
\end{small}

\begin{theorem} The construction in Figure\ \ref{fig_folding} is a complete and 
knowledge sound IVC scheme.
\end{theorem}

The completeness of the construction is apparent. As the foldpot scheme
is an integration of SuperNova (extension of Nova) and Cyclefold,
its soundness proof originates from Lemma 10 of Nova \cite{Nova}, 
Lemma 3 of SuperNova \cite{SuperNova}, and Lemma 
2 of CycleFold \cite{CycleFold}, with adaptation from CCS to standard
R1CS. For the soundness proof, the key observation is that given IVC proof  
$\ivcprf_{i} = \left(
	\begin{array}{l}
		\left(\U_{i}, \W_{i}, \u_{i}, \w_{i}\right), 
		\left(\Uc{i}, \Wc{i}\right), 
		\left(\Upair{i}, \Wpair{i}\right),
		\pc_{i} 
	\end{array}
\right)$, one can extract $\U_{i-1}$ and $\u_{i-1}$ which
fold into $\U_{i}$. Then based on the extractor properties for
folding relations, one can extract the $\W_{i-1}$ and 
$\w_{i-1}$ for $\U_{i-1}$ and $\u_{i-1}$. The same observation
applies to the cyclefold and cycle-pair relations. Thus,
it allows the knowledge extraction to proceed recursively.
The cyclefold instance, according to Lemma 2 of CycleFold \cite{CycleFold}
guarantees the entire scheme's knowledge extraction, by asserting
the correctness of the out-sourced linear combination of curve points.
The correctness of cyclepair is done through the check of
$\z_i = \hash(\O_i)$ in the $\augf{j}$ circuit.

\subsection{Circuit}
In this section we provide further details about the circuits.
Given a set of words $\myvec{s} = \myvec{s}^{(1)} \concat ... \concat \myvec{s}^{(n)}$, our goal is to prove that a certain predicate $\pred{\myvec{s}}$ 
are satisfied by these words.  As an example, the predicate may assert that
each word does not match any of the regex in a malware signature set.
As another example, the predicate asserts that
the words correctly encode the verifier function of a QA-NIZK relation.
We accomplish the ``proof" of $\pred{\myvec{s}}$ via
a sequential execution of a collection of circuits
$\{\circuit_1, ... \circuit_k\}$ via $N$ steps s.t.
$\forall i\in [N]: z_i = \circuit^i_{\pc_i}(z_{i-1})$,
and in the last output $z_N$ there is a boolean flag
indicating if $\pred{\myvec{s}}$ is true. We use IVC here for
the folding speed is faster than general purpose SNARK provers.
However, the computing model is closer to the Repeated Computation
with Global State (RCG) \cite{RCG} in the way that 
the prover witnesses of all the steps are already known before folding 
is started, and so is zk-rollup \cite{Rollup}.

\subsection{Circuit and Function View of $\Sigma$-IR1CS}
It is convenient to view a $\Sigma$-IR1CS in the point of view of
functions and also as an arimetic circuit. 
We mainly consider functions (circuits, i.e., $\Sigma$-IR1CS)
that can compute values given input words. To simplify 
the presentation, we assume that these functions process
fixed size words, e.g., a circuit that discharge a word of
size $2^{10}$ of a given set of malware signatures.\footnote{In 
implementation, we allow circuit to process a word with a prior-set
max length, so that we can reduce the number of circuits needed
to handle words of arbitrary length, i.e., we do not have to
provide circuits which handles word length $2^0$, $2^1$, etc
which are not ``economical" given the IVC overhead.} 
Note that however, the word problem itself, is not fixed size.
We thus distinguish these two cases in the following definition.

\begin{definition} A {\em word processing function} 
has the form
$z \leftarrow f(\myvec{w}, \myvec{p})$, with linear time and
space copmlexity.  Here $z\in \F$, $\myvec{w} \in \F^n$ and
$\myvec{p} \in \F^m$ represent the output, the input word,
and any additional information. If the function 
needs to return multiple elements, we let a fixed structure
$\ds{z}$ encode them s.t.  $z = \hash(\ds{z})$. We call
$\ds{z}$ its {\em output structure}.
$f$ is called of {\em fixed size ($n$)} if $n$ is a fixed constant. 
\end{definition}

Since our goal is RCG (IVC), we would like to split an expensive
function into many steps.  We say that a word function 
$f(\myvec{w}, \myvec{p})$ is {\em streaming}, if it can
process the word ``character by character". A streaming
function is a stronger notion than recompusively computable by IVC,
in fact, it implies that the function
can be computed chunk by chunk (more efficient than character by character).
Streaming functions enjoy a nice property: 
combination of any two streaming functions is still
streaming.\footnote{Note that this property does not apply to
any recursively computing functions. If two functions each has
special requirement on the chunk size, then they cannot be safely composed.}

\begin{definition} A word function $f$
is called {\em streaming}, if there exists 
a function $g$ of fixed-size $1$,
whose output structure contains an element $v$
s.t. for any $\myvec{w}, \myvec{p}$, letting $n = |\myvec{w}|$,
there exists a sequence of $z_0, ..., \z_{n}$ satisfying:
\begin{enumerate}
	\item$\forall i in [1,n]: z_i = g(\myvec{w}_i, \myvec{p})$.
	\item $\ds{z_{n}}.v = f(\myvec{w}, \myvec{p})$.
\end{enumerate}
We call $g$ the atomic function of $f$.
\end{definition}

\begin{lemma} Given any two streaming functions $f(\myvec{w}, \myvec{p})$
and $g(\myvec{w}, \myvec{p})$, their {\em composition} 
(denoted as $f \circ g$) is defined
as a function $z\leftarrow h(\myvec{w}, \myvec{p})$ s.t. $\ds{z}$
is a vector $\big(f(\myvec{w}, \myvec{p}), g(\myvec{w}, \myvec{p})\big)$.
$h$ is streaming.
\end{lemma}

All the approximated regex discharging algorithms are streaming functions.
Streaming implies recursive computation using IVC.
When processing a word problem, there is a need to pass
information between circuits, e.g., the remaining signatures
that are not handled by $\cp$ will be passed to $\bw$, and then 
$\sde$ circuits. 
\footnote{Notice that, these input/output information belong to the
{\em fixed meomry} part of the non-deterministic advice
to the circuit, i.e., they do not depend on verifier challenges.}
We thus further (synatically) restrict $\Sigma$-IR1CS relations
as special arithmetic circuit with defined structures.

\begin{definition}  Given a fixed-size streaming function $f$,
we say a $\Sigma-$ I-R1CS instance 
$(N, n, \F; \myvec{T} \in \F^{N\times 2}, $
$\ell_x, \ell_{m_1}, \ell_{m_2}, \ell_{m_3} \in \nat; 
A,B,C \in \F^{n\times m}; \chi)$
 is the arithmetic circuit for $f$, denoted as $\circuit_{f}$,
if $\ell_x = \ell_w + \ell_i + \ell_o + \ell_d + 2$, i.e.,
its statement $\myvec{x}$ is structured as the concatenation
of $4$ (fixed size) segements $\xw$ (the word),
$\xi$ (the input), $\xo$ (the output), $\xd$ (the data),
and two additional field elements $z_{i-1}$ and $z_{i}$,
and for any satisfying instance for $\circuit_{f}$, its
$z_i = f(z_{i-1}, \xw \concat \xi \concat \xo \concat \xd)$. 
\end{definition}

Intuitively, given a streaming function $f$, and its atomic
function $g$. Let $g^{(i)}$ be a function by chaining $g$ for $i$ times.
We have a collection of circuits 
$\left\{\circuit_{g^{(2^i)}}\right\}_{i=1}^k$
that can compute $f$ for any input word. A larger $k$ permits
less overhead caused by IVC in practice.


\subsection{Batch and Individual Proof}
We recall the idea of batch and individual proof and formally define the
proof framework in this section. There is a public regex set
$R$, whose approximated discharging algorithms can be encoded
as a streaming function which utilizes
a lookup table with two columns $\lkup = (\lk{1}, \lk{2})$, 
denoted as $f_\lkup$. The prover holds
a {\em secret} set of words 
$\left\{\word{i} \right\}_{i=1}^n$, and alternatively
we write it as $(\fullword, \veclen)$ where 
$\fullword = \word{1} \concat ... \concat \word{n}$ and
$\forall i\in[n]: \veclen_i = |\word{i}|$. 
The proof proceeds in two phases: (1) {\em batch proof:} where
the prover produces one single proof for $\fullword$ that each of its
$\word{i} \not\in R$. (2) {\em individual proof:}
where given $\bargen{\word{i}}$, the prover convinces the verifier that
it indeed encodes is the i'th word in $\bargen{\fullword}$.
This scheme is accomplished using a $O(1)$ size SNARK proof (decider
proof for an IVC scheme) and several KZG proofs.

Now consider the individual proof for $i$. The verifier sees:
$\bargen{\word{i}}$, $\bargen{\fullword}$, $\bargen{\veclen}$.
The prover can produce a non-interactive proof as below:
let $r_i = \rohash\left(i, \bargen{\word{i}}, \bargen{\fullword}, 
\bargen{\veclen}\right)$.
let $v_i = \sum_{j=0}^{|\word{i}|-1} r^i \word{i}_{j}$.
The prover can provide a kzg proof to show that
$\bargen{\word{i}}$ evaluates to $v_i$ at $r_i$.
Now it remains to show that according to 
$\veclen$, the corresponding $i$'th segment
of $\fullword$ is indeed $\word{i}$.
This is accomplished combined with the SNARK of the 
batch proof. In the batch proof, there is an IVC sequence
of circuits (whose concatenation of $\xw$ is $\fullword$).
We perform in circuit evaluation of $r_i$ for each $\word{i}$.
Thus, the verification of individual proof relies on the batch proof,
which we elaborate the details below.

We first formally define the statement of a batch proof.
We say a streaming predicate is a streaming function which 
returns a booleaning value.
Fix a streaming predicate $\pred{\lkup}$, 
given any word problem $(\fullword, \veclen)$,
we define the public batch proof claim 
$\pubstat_{(\fullword, \veclen, f_\lkup)}$
(shortened by $\pubstat$ when the context is clear) 
as a tuple:
$
	(\bargen{\fullword}, \bargen{\veclen}, \bargen{\lkup},
		\bargen{\myvec{r}}, \bargen{\myvec{v}})
$, where $\myvec{r}$ and $\myvec{v}$  encodes the
$r_i$ and $v_i$ for the individual proof of each word.
$\pubstat$ is visible to verifier and all commitments are completely
hiding KZG commitments. The relation $R_\pubstat$ is defined
as below:
	%(\bargen{\fullword}, \bargen{\veclen}, \bargen{\lkup},
	%	\bargen{\myvec{r}}, \bargen{\myvec{v}})
\begin{equation}
	R_{\pubstat_{\fullword, \veclen, \pred{\lkup}}} = 
	\left\{ 
		(\bargen{\fullword}
		\bargen{\veclen}, 
		\bargen{\lkup},
		\bargen{\myvec{r}}
		\bargen{\myvec{v}}) \middle\vert
			\begin{array}{l}
				\fullword = \word{1} \concat ... \concat 
					\word{(|\veclen|)} \\
				\bargen{\fullword} = \kzgcom(\pk, \fullword) \\
				\bargen{\veclen} = \kzgcom(\pk, \veclen) \\
				\bargen{\myvec{r}} = \kzgcom(\pk, \myvec{r}) \\
				\bargen{\myvec{v}} = \kzgcom(\pk, \myvec{v}) \\
				\forall i\in [|\veclen|]: \veclen_i = |\word{i}| \\
				\forall i\in [|\veclen|]: \myvec{v_i} = \sum_{j=0}^{|\veclen_i|}
					(\myvec{r}_i)^j \cdot \word{i}_j \\
				\forall i\in [|\veclen|]: \exists \myvec{p}:
					\pred{\lkup}(\word{i}, \myvec{p}) = 1
			\end{array}
	\right\}
\end{equation}

Intuitively, $R_\pubstat$ states that each word satisfies
the predicate, the vector lengths match the words,
and given a non-deterministic vector $\myvec{r}$, the
weighted sum of each $\word{i}$ using $\myvec{r}_i$
generates $\myvec{v}_i$.\footnote{We do not have to
specify that $r_i$ is the result of hashing individual claim.
This can be checked in the individual proof using vector commitments.}
It is not hard to see that enforcing the $(\myvec{r}, \myvec{v})$ semantics
itself is a streaming function (essentially computing weighted sum).
Combining all required checks (including $\pred{\pubstat}$) results
in a streaming function, which can be recursively computed through IVC.
This immediately applies to the approxmiated Regex discharging problem.

\begin{theorem} \label{thm:main}
Given any streaming predicate $\pred{\lkup}$, there exists 
a set of circuits $\myvec{\circuit}$ and an IVC scheme, 
such that for any word problem
$(\fullword, \veclen)$, the IVC scheme using $\myvec{\circuit}$
can produce a zk-SNARK proof 
for $R_{\pubstat_{\fullword, \veclen, \pred{\lkup}}}$.
The the prover work is $O(|\fullword|)$ group operations, 
the proof size and verifier time are $O(1)$. 
\end{theorem}

We now introduce the an IVC scheme that realizes Theorem\ \ref{thm:main},
and present the scheme in Figure\ \ref{fig:main}. It is a concrete
scheme that realizes Theorem\ \ref{thm:main}. It has four stages.

In Stage 1: given a collection of streaming circuits $\{\circuit_i\}_{i=1}^k$
which computes $f_{R_\pubstat}$, we run the standard RCG scheme. 
At step 4, we first run a heurstics, depending on the $\cp$, $\sde$, $\dfa$
proofs for each word and the capacity of circuits, determine the 
$\pc_i$ for each step to minimize resource consumption. These discharging
proofs are encoded as fixed witness $\myvec{w}^{(\rho,i)}$, and at Step 6
we compute $\hcf$ as the hashchain of the fixed segments first for
Fiat-Shamir randoms.  Note that we also embed into the fixed segments:
(1) each word $\word{i}$, (2) its length $\veclen_i$, and (3)
the randoms and evals of $\myvec{r}_i$ and $\myvec{v}_i$ used for
individual proof. Then the random challenge $r$ and a combination 
factor $c$ is generated over $\hcf$. We run the IVC scheme, and after
$N$ steps, at step (8), the ivc proof
$\ivcprf_N$ asserts the validity of a variety of attributes
about the words such as $\veclen{i} = |\word{i}|$, and
$\myvec{v}_i = \sum_{j=0}^{|word{i}|-1} \myvec{r_i}^{j-1} \cdot \word{i}[j]$.
Note $\ds{z_N}$ also includes attribute $e$ as the
evaluation of a combination of vectors (polynomials), which 
we later correlate with KZG proof.

Now at the end of Stage 1, we have a SuperNova IVC proof, which has
$n$ instances of committed R1CS, where each has a
$\overline{\myvec{W}^{(i)}}$,
$\overline{\myvec{E}^{(i)}}$,
$\overline{\myvec{W}^{(\rho, i)}}$. We need to demonstrate the knowledge
of the correpsonding witness for these commitments. In
Sonobe \cite{Sonobe}, this is handled by providing one KZG commitment
for each witness, and do a matching evaluation in circuit. This works
when $k$ is small, but will result in a large proof size (e.g., $60$ KZG
proofs for $20$ circuits!)

Stage 2: We take an alternative approach 
using QA-NIZK proof. The basic idea is that
we collect the concentation of all witnesses and call it
$\myvec{W} \concat \myvec{E}$ (see Step 11). We then provide
a matrix based QA-NIZK proof which shows that the corresponding
subsegments of $\myvec{W} \concat \myvec{E}$ generates 
all the $3k$ commitments. The verification equation takes
$3k+1$ pairing operations (see Step 12). 

We would like to further compress the proof and verification cost 
of Step 12. This is done through Stage 2 (steps 11 to 12) where
we run another $3k+1$ steps of IVC, using one single chip
for $\cyclepair$, to prove the QA-NIZK equations. Essentially it is
to out-source one pairing operation to Grumpkin curve in each step to
minimize folding cost. Apparently, when $k$ is large, the overall
prover cost is much smaller compared with encoding all $3k+1$
non-natively over EC1 (e.g., BN254).

Stage 3: We now have two IVC proofs: $\ivcprf_N$ for stage 1 (main circuits),
and $\ivcprf_k$ for the cyclepair circuit. We then generate a SNARK
proof $\prfsnark$ (Lines 15-18), which proves the knowledge
of these two pairs of IVC proofs, that satisfy a certain properties:
(1) $\ds{z_N}.e$ matches the given $e$, which is the combined
evaluation of polynomials behind $\pubstat$, and (2)
$\ds{z_N}.{\pred{\pubstat}}$ is true.

Recall that the decider of IVC mainly performs two tasks: (1) to prove
that the witness satisfy the circuit relations, and (2) to prove that
the witness does generate the commitments in the committed SuperNova instances.
Here we let $\rhalf$ to represent the logic of (1), and use the
Sonobe technique to outsource (2) to KZG commitment. The idea is:
let $\overline{\myvec{W} || \myvec{E}}$ be the KZG commitment to
all witness and error terms of all circuits. We provide a KZG proof
for showing that it evaluates to a certain value $e_2$ at a point $r_2$.
Then in circuit, we evaluate the polynomials and prove that all witness/error
terms do evaluate to $e_2$ at $r_2$. We include such KZG proof
in $\prfsnark$ at line (14), and similarly do this for the 
cyclepair circuit using $(r', e')$.

Stage 4: we build the batch proof as planned, which centeres around
the $\prfsnark$, and provide the corresponding KZG proof.
As the 2nd stage SuperNova has only one circuit, we disclose its
hiding commitments, and use an extra QA-NIZK proof to prove
the knowledge of witness behind these commitments with
$\overline{\myvec{W'} \concat \myvec{E'}}$.
For individual proof, we mainly re-generating the public randoms,
$\myvec{r}_i$. Show that the committed word $\overline{\word{i}}$
does generate $\myvec{v}_i$, and show that they are the
$i'$th pair in the vaues embedded in the circuit, which links
$\word{i}$ with the $i$'th word in the circuit.

The batch proof size is: 13 $\G_1$, 1 $\G_2$ and 5 $\F$ elements.
Verifier cost is mainly 10 pairings.
Individual proof size is 2 $\G_1$, and 1 $\F$. Verifier cost
is: 4 pairings.


\def\hp{\hspace*{0.1in}}
\begin{scriptsize}
\begin{algorithm*}[bt]
\SetKwData{Left}{left}\SetKwData{This}{this}\SetKwData{Up}{up}
\SetKwFunction{Union}{Union}\SetKwFunction{FindCompress}{FindCompress}
$ \left(\big(\pubstat, \prfbatch\big), \left(\pubstatind{i},\prfind{i}\right)_{i=1}^{n}\right) 
	{\leftarrow} \mathsf{prove\_all}\left(\pk, 
		\big\{\word{i}\big\}_{i=1}^n, \lkup, \left(\circuit_i\right)_{i=1}^k 
		\right)$ 
\newline
(1) Compute $\veclen \leftarrow \left(|\word{i}|\right)_{i=1}^n$, and
	$\fullword \leftarrow \word{1} \concat ... \concat \word{n}$. \newline
(2) Generate KZG commitments: 
	$\left(\bargen{\word{i}}\right)_{i=1}^n$, 
	$\bargen{\veclen}$, 
	$\bargen{\fullword}$, $\bargen{\lkup}$.
	Compute $\myvec{r} \leftarrow \left(\rohash\big(i, \bargen{\word{i}}, \bargen{\fullword}, \bargen{\veclen}\big)\right)_{i=1}^n$, 
	and 
	$\myvec{v} \leftarrow \left(\sum_{j=0}^{\veclen_{i}-1} \myvec{r}_i^j \word{i}_j\right)_{i=1}^n$. \newline
(3) Compute KZG for $\myvec{r}$ and $\myvec{v}$.
	Define $\pubstat = (\bargen{\fullword}, \bargen{\veclen}, \bargen{\lkup},
		\bargen{\myvec{r}}, \bargen{\myvec{v}})$,
	and for $i\in [n]:$ $\pubstat_i = \left(\bargen{\fullword}, \bargen{\veclen}, \bargen{\word{i}}\right)$. \newline
\# {\bf    Stage 1: main proof} \newline
(4) Based on capacity of $\left(\circuit_i\right)_{i=1}^k$,
	run heurstics to determine $N$, the steps needed for
	computing $f_{R_{\pubstat}}$, and 
	$\left(\pc_i\right)_{i=1}^N$, the circuit IDs of all steps 
	\newline
(5) For $i\in [n]$: Run $\cp$, $\sde$, $\dfa$, $\isde$ for each $\word{i}$,
	and generate non-deterministic advice $\omega^{(i)}$ for
	each $f_{\circuit_{\pc_i}}$, and based on it,
	generate the partial witness $\myvec{W}^{(\rho,i)}$,
	includes $\word{i}$, $\myvec{r}_i$, and $\myvec{v}_i$,
	and $\lkup^{(i)}$ (the $i$'th piece of lookup table distributed
	into $\circuit_{\pc_i}$), $\qry{i}$ (the query table 
	in circuit $\circuit_{\pc_i}$). We write
	$\fullqry = \qry{1} \concat ... \concat \qry{N}$,
	and note $\fullqry \lookup \lkup$.
	\newline
(6) Compute $\hc \leftarrow \hashchain(\{\myvec{W^{(\rho,i)}}\}_{i=1^N})$.
Then $(r,c,\beta) \leftarrow \rohash(\hc)$. \newline
(7) Compute $z_0 \leftarrow \hash(\ds{\z_0})$, where
$\ds{\z_0}$ contains constants $(\hc, r, c, \beta)$ that
are passed to each $\ds{z_i}$ and variables for computing
accumulated sum such as $\sum_{i=1}^N \fullword_i r^{i-1}$. 
Let $\ivcprf_0$ be trivially satisfying IVC proof for initialization.
\newline
(8) For $i\in [N]:$ compute the full non-deterministic advice 
$\omega^{(i)}$ with $(\hc, r,c,\beta)$, and $z_{i+1} \leftarrow f_{R_\pubstat}\left(z_i, \omega^{(i)}\right)$. Run $\ivcprf_{i+1} \leftarrow \ProveIVC\left(\pk, (i, z_0, z_i), \omega^{(i)}, \pc_{i+1}, \ivcprf_i\right)$. 
A valid $\ivcprf_N$ proves that via $N$ steps if IVC, 
$z_N \leftarrow f_{R_\pubstat}(\fullword, \veclen)$, 
where $b_{\pubstat}$ in $\ds{z_N}$ asserts
$
 \left(
 \begin{array}{l}
 	\forall i\in[n]: \veclen_i = |\word{i}| \\
 	\forall i\in[n]: \myvec{v}_i = \sum_{j=1}^{|\word{i}|} 
		\word{i}_j (\myvec{r}_i)^{j-1} \\
	\sum_{i=1}^{|\fullqry|} \frac{1}{\beta + \fullqry^{(1)}_i + \alpha \cdot \fullqry^{(2)}_i} = 
	\sum_{i=1}^{|\lkup|} \frac{\myvec{m}_i}{\beta + \lkup^{(1)}_i + \alpha\cdot \lkup^{(2)}_i} \\
	\forall i\in[n]: \cp(\word{i}) \subseteq \sde(\word{i}) \union \dfa(\word{i}) \union \isde(\word{i}) 
 \end{array}
 \right)
$
and $v$ in $\ds{z_N}$ has value 
$\left(\begin{array}{l}
c^0 \cdot \sum_{i=1}^{|\lkup|} \lkup^{(1)}_i \cdot r^{i-1} \\ 
+ c^1 \cdot \sum_{i=1}^{|\lkup|} \lkup^{(2)}_i \cdot r^{i-1} \\ 
+ c^2 \cdot \sum_{i=1}^{|\fullword|} \fullword_i \cdot r^{i-1} \\ 
+ c^3 \cdot \sum_{i=1}^{|\veclen|} \veclen_i \cdot r^{i-1}\\
+ c^4 \cdot \sum_{i=1}^{|\veclen|} \myvec{r}_i \cdot r^{i-1}\\
+ c^5 \cdot \sum_{i=1}^{|\veclen|} \myvec{v}_i \cdot r^{i-1}
\end{array}
\right)
$
\newline
\# {\bf    Stage 2: CyclePair Proof For Saving Verifier Cost} \newline
(10) Parse $\ivcprf_N$ as $\left(
		\left(\U_N, \W_N, \u_N, \w_N\right), 
		\left(\Uc{i}, \Wc{i}\right), 
		\left(\Upair{i}, \Wpair{i}\right),
		\pc_N, \O_N
	\right)$. 
	Compute $(\U_{N+1}, \W_{N+1})$ by folding
	$(\U_N, \W_N)$ and $(\u_N, \w_N)$. 
	For $j\in [k]$: extract the comitments: 
		$\left(\overline{\myvec{W}^{(j)}},
		\overline{\myvec{E}^{(j)}},
		\overline{\myvec{W}^{(\rho,j)}}\right)$ from $\U_{N+1}[j]$
		and witness and error terms 
		$\left(\myvec{W}^{(j)}, \myvec{E}^{(j)}\right)$
		from $\W_{N+1}[j]$.
	Let $\myvec{W} = \myvec{W}^{(1)} \concat ... \concat \myvec{W}^{(k)}$, i.e.,
	concatenated witness. Note that
	$\myvec{W}^{(\rho,i)}$ is a subsegment of $\myvec{W}^{(i)}$ located
	at fixed position.
	Similarly let 
	$\myvec{E} = \myvec{E}^{(1)} \concat ... \concat \myvec{E}^{(k)}$. 
	Compute $\overline{\myvec{W} \concat \myvec{E}}$ as a hiding KZG
	commitment of $\myvec{W} \concat \myvec{E}$.
	\newline
(11)  Construct QA-NIZK Proof $\prfqa$ for the following relation that
	asserts the KZG commitment $\overline{\myvec{W} \concat \myvec{E}}$
	with all commitments in supernova commmitted instance $\U_{N+1}$:
	$\left\{    
		\left((\myvec{W}^{(j)}, \myvec{E}^{(j)})_{j=1}^k\right):
		\forall j\in [k]:
			\left(
			\begin{array}{l}
			\overline{\myvec{W}^{(j)}} = \kzgcom(\myvec{W}^{(j)}). \\
			\overline{\myvec{W}^{(\rho,j)}} = \kzgcom(\myvec{W}^{(\rho,j)}). \\
			\overline{\myvec{E}^{(j)}} = \kzgcom(\myvec{E}^{(j)}),
			\end{array}
			\right), \textrm{ and }
			\overline{\myvec{W}\concat\myvec{E}} = \kzgcom(
				\myvec{W}^{(1)} \concat ... \concat \myvec{W}^{(k)} \concat
				\myvec{E}^{(1)} \concat ... \concat \myvec{E}^{(k)}
			).
	\right\}
	$, where
	its public input $\xqa$ is: 
		$\left(\overline{\myvec{w} \concat \myvec{E}},
		\left(\overline{\myvec{W}^{(j)}},
		\overline{\myvec{E}^{(j)}},
		\overline{\myvec{W}^{(\rho,j)}}\right)_{j=1}^k
		\right)$.
	Note that its verifier key $\vkqa = (\myvec{C}, a)$ 
	is a $3k+2$ vector of $\G_2$ elements, with $|\myvec{C}| = 3k+1$.
	We now construct another IVC to prove the QA-NIZK equation
	$\sum_{j=1}^{3k+1} e(\xqa_j, \myvec{C}_j) = e(\prfqa, a)$. 
\newline
(12) Repeat step (4) to (8) for a single circuit system
	where the circuit encodes $\cyclepair$. Running $k$ steps
	of IVC, results in the assertion for public
	statement: $(\hc_1, \hc_2)$, where there exists
	$\xqa\in \G_1^{3k+1}$ and $\myvec{C} \in \G_2^{3k+1}$ s.t.
	$\sum_{j=1}^{3k+1} e(\xqa_j, \myvec{C}_j) = e(\prfqa, a)$, i.e.,
	the QA-NIZK verification equation. \newline
{\bf \# Stage 3. Generate SNARK Proof} \newline
(13) Given a supernova IVC proof $\ivcprf$ for a scheme with $k$
instances, let predicate $\rhalf(\ivcprf_i)$ defines the ``half" 
of the SuperNova decider check on $\ivcprf_i$ by skipping
the check on witness commitments. 
Concretely it is define as as below:
Parse $\ivcprf_i = \left(
		\left(\U_i, \W_i, \u_i, \w_i\right), 
		\left(\Uc{i}, \Wc{i}\right), 
		\left(\Upair{i}, \Wpair{i}\right),
		\pc_i, \O_i
	\right)$. 
Compute (non-natively for group operations in circuit)
$(\U_{i+1}, \W_{i+1}, \u_{i+1}, \w_{i+1})$ by folding
$(\U_i, \W_i, \u_i, \w_i)$. 
Use non-deterministic advice $\ds{z_i}$ to show
to show:  (1) $(\u_i.\overline{E}, \u_i.u) = (\u_\bot.\overline{E}, 1)$.
(2) $\u_{i+1}.x = \left(\hash(\vk, i, z_0, z_i, U_{i+1}), 
\hash(\vk, \Uc{i}), \hash(\vk, \Upair{i}) \right)$.
(3) $\hash(\ds{z_i}) = z_i$.
(4) $\forall j\in [k]:$ $\U_{i+1}[j]$ and $\W_{i+1}[j]$ is
	a relaxed R1CS instance for circuit $\circuit_{j}$.
(4) $\Uc{i}$ and $\Wc{i}$ is a relaxed cyclefold instance.
(5) $\Upair{i}$ and $\Wpair{i}$ is a relaxed cyclepair instance. \newline
 \hspace*{0.1in} Now given $\ivcprf_N$ (the IVC proof of stage 1),
 	and given $\myvec{W} \concat \myvec{E}$ in Step 10.  Let
	$\ds{z_N}$ be the non-deterministic advice for $z_N = \hash(\ds{z_N})$.
	Compute $r_2 \leftarrow 
		\rohash(\ds{z_N}.\overline{\myvec{W}^{(\rho)}}, 
		\overline{\myvec{W} \concat \myvec{E}})$.
	Compute $e_2 \leftarrow \sum_{i=1}^{|\myvec{W} \concat \myvec{E}|}
		(\myvec{W} \concat \myvec{E})_i \cdot r_2^{i-1}$. 
	Compute the correpsonding KZG proof $\prfkzgwe$ for $e_2$.
	Similarly for $\ivcprf_k'$ (the IVC proof for stage 2),
	given $\ds{z_k'}$, $\myvec{W'} \concat \myvec{E'}$ extracted
	from $\ivcprf_k'$,
	compute $r' \leftarrow \rohash(\ds{z_k}'.\overline{\myvec{W'}^{(\rho)}}
		, \overline{\myvec{W'} \concat \myvec{E'}})$, and
		$e' \leftarrow \sum_{i=1}^(\myvec{W'} \concat \myvec{E'})_i \cdot
			{r'}^{i-1}$.
		Similarly compute the $\prfkzgwetwo$ for
		$e'$.
	Let $(\U_{i+1}', \W_{i+1}')$ be the folded instance for $\ivcprf_{k}'$
	as shown in computing $\rhalf$.
	Since stage (2) has one circuit, collect
		$(\overline{\myvec{W'}}, \overline{\myvec{E'}}, 
			\overline{\myvec{W'}^{(\rho)}})$ from its own
			$\W_{i+1}'[0]$. 
	Similarly to step (11), construct $\prfqawetwo$
	for showing the relation between 
	$\overline{\myvec{W'} \concat \myvec{E'}}$,
	and $\overline{\myvec{W'}}$,
	$\overline{\myvec{E'}}$,
	$\overline{\myvec{W'}^{(\rho)}}$.
	Now construct the public
		statement for the SNARK proof as below:
	$
		\statsnark = \left( 
			\left(
				(r,c,\ds{z_N}.v), (r_2,e_2), \ds_{z_N}. \overline{\myvec{W}^{(\rho)}}
			\right),
			\left(
				(r',e'), (\overline{\myvec{W'}}, \overline{\myvec{E'}},
					\overline{\myvec{W'}^{(\rho)}}
			\right)
		\right)
	$ \newline
(14) Produce $\prfsnark$ for $\statsnark = 
		\statsnark = \left( 
			\left(
				(r,c,e), (r_2,e_2), \overline{\myvec{W}^{(\rho)}}
			\right),
			\left(
				(r',e'), (\overline{\myvec{W'}}, \overline{\myvec{E'}},
					\overline{\myvec{W'}^{(\rho)}}
			\right)
		\right)$,
i.e., there exists
$\ivcprf_N$ and $\ivcprf_k'$ which satisfies the following:
(1) $\rhalf(\ivcprf_N)$ for stage 1, (2) $\rhalf(\ivcprf_k')$ for 
stage 2, (3) $(r,c,e, \overline{\myvec{W}^{(\rho)}}) \in \ds{z_N}$, 
(4) Given $r_2$, and $\myvec{W} \concat \myvec{E}$ collected from $u_N$, verify
$e_2$ is equal to the value computed in Step (13).
(5) Similarly verify $e'$ using formula given in Step (13).
(6) Assert $\overline{\myvec{W}'}, \overline{\myvec{E'}},
	\overline{\myvec{W'}^{(\rho)}} = \u_k'[0]$  \newline
\# {\bf Stage 4: Build Batch and Individual Proof} \newline
(15) Compute 
	$\barvec{\myvec{r}} \leftarrow \veccom(\myvec{r})$,
	$\barvec{\myvec{v}} \leftarrow \veccom(\myvec{v})$.
	Compute QA-NIZK proof $\prfqarv \in \G_1$ for
	showing their equivalence to KZG, i.e., there exist $\myvec{r}$
	which generates both $\barvec{\myvec{r}}$ and $\bargen{\myvec{r}}$ (KZG),
	and also $\myvec{v}$
	which generates both 
	for $\barvec{\myvec{v}}$ and $\bargen{\myvec{v}}$. \newline
(16) Comute batched KZG proof $\prfkzgagg$, which asserts
that given $\pubstat = \left(\overline{\lkup^{(1)}}, 
	\overline{\lkup^{(2)}}, \overline{\fullword}, 
	\overline{\veclen}, \overline{\myvec{r}},
	\overline{\myvec{v}} \right)$, combining the evaluation
	at $r$ using combinning factor $c$, results in
	$e$, Formally, let $\myvec{p}$ be the corresponding
	polynoaims for $\pubstat$, $e = \sum_{i=1} c^{i-1} \cdot
	\myvec{p}_i(r)$,i.e., it matches the $\ds{z_N}.v$
	in Step (9). \newline
(17) Output $\prfbatch = \left(
     \hc(\{\overline{\myvec{W}^{(\rho,i)}}\}_{i=1}^N),
	 \left(\barvec{\myvec{r}}, \barvec{\myvec{v}}, \prfqarv\right),
	 \left(r,c,e,\prfkzgagg\right),
	 \left(\overline{\myvec{W} \concat \myvec{E}}, r_2, e_2,
	 	\prfkzgwe \right),
	 \left(\overline{\myvec{W'} \concat \myvec{E'}}, r', e',
	 	\prfkzgwetwo \right),
	 \left(\overline{\myvec{W'}}, \overline{\myvec{E'}}, 
	 	\overline{\myvec{W'}^{(\rho)}},
		\prfqawetwo
	 \right),
	 \prfsnark
	\right)$. \newline
(18) For each $i\in [n]:$ 
	Compute $c_i \leftarrow \rohash(\myvec{r}_i, \myvec{v}_i)$, 
	$\prfveccom^{(i)} \leftarrow \veccomprove(\pk, 
		\myvec{r} + c_i \cdot \myvec{v}, 
		i, 
		\myvec{r}_i + c_i\cdot \myvec{v}_i)$.
	$\prfkzgwi \leftarrow \kzgcomprove(\pk, \myvec{w}^{(i)}, \myvec{r}_i,
		\myvec{v}_i)$.
	Output $\prfind{i} = \left(\prfveccom^{(i)}, \myvec{v}_i, \prfkzgwi
	\right)$.	

$ 1/0 {\leftarrow} \mathsf{verify\_batch}\left(\vk, 
		\pubstat,
		\prfbatch	
		\right)$ 
\newline
(1) Parse $\pubstat = \left(\overline{\lkup^{(1)}}, 
	\overline{\lkup^{(2)}}, \overline{\fullword}, 
	\overline{\veclen}, \overline{\myvec{r}},
	\overline{\myvec{v}} \right)$, 
	and 
	$\prfbatch = \left(
     \hcf,
	 \left(\barvec{\myvec{r}}, \barvec{\myvec{v}}, \prfqarv\right),
	 \left(e,\prfkzgagg\right),
	 \left(\overline{\myvec{W} \concat \myvec{E}}, e_2,
	 	\prfkzgwe \right),
	 \left(\overline{\myvec{W'} \concat \myvec{E'}}, e',
	 	\prfkzgwetwo \right),
	 \left(\overline{\myvec{W'}}, \overline{\myvec{E'}}, 
	 	\overline{\myvec{W'}^{(\rho)}},
		\prfqawetwo
	 \right),
	 \prfsnark
	\right)$. \newline
(2) Compute $(r,c) \leftarrow \rohash\left(\pubstat, \hcf\right)$,
	and $r_2 \leftarrow 
		\rohash\left(\hcf, 
		\overline{\myvec{W} \concat \myvec{E}}\right)$, and
	$r' \leftarrow \rohash\left(\overline{\myvec{W'}^{(\rho)}}
		, \overline{\myvec{W'} \concat \myvec{E'}}\right)$.  \newline
(3) Assert $\qaverify\left(\vk, 
		\left(\barvec{\myvec{r}}, \barvec{\myvec{v}},
	\bargen{\myvec{r}}, \bargen{\myvec{v}}\right), \prfqarv\right)$.
 		and $\qaverify\left(\vk, 
	 		\left(\overline{\myvec{W} \concat \myvec{E}}, 
	 			\overline{\myvec{W'}}, \overline{\myvec{E'}}, 
	 			\overline{\myvec{W'}^{(\rho)}}
			\right), \prfqawetwo\right)$. \newline
(4) Assert $\kzgcomverify(\vk, \sum_{i=1}^6 \pubstat_i \cdot c^{i-1}, r, e, \prfkzgagg)$, and 
	$\kzgcomverify(\vk, \overline{\myvec{W} \concat \myvec{E}}, r_2, e_2, 
		\prfkzgwe)$, and
	$\kzgcomverify(\vk, \overline{\myvec{W'} \concat \myvec{E'}}, r', e', 
		\prfkzgwetwo)$. \newline 
(5) Construct $
		\statsnark = \left( 
			\left(
				(r,c,e), (r_2,e_2), \overline{\myvec{W}^{(\rho)}}
			\right),
			\left(
				(r',e'), (\overline{\myvec{W'}}, \overline{\myvec{E'}},
					\overline{\myvec{W'}^{(\rho)}}
			\right)
		\right)$, Assert $\snarkverify(\vk, \statsnark, \prfsnark)$.

$ 1/0 {\leftarrow} \mathsf{verify\_individual}\left(\vk, 
		\pubstat,
		\pubstatind{i},
		\prfind{i},
		\right)$ 
\newline
(1) Parse $\pubstatind{i} = \left(i, \overline{\word{i}}, \overline{\fullword},
	\overline{\veclen}\right)$, 
	and 
	$\prfind{i} = \left(\prfveccom^{(i)}, \myvec{v}_i, \prfkzgwi
	\right)$.	 \newline
(2) Compute $\myvec{r}_i \leftarrow \rohash(\pubstatind{i})$.
Assert $\kzgcomverify(\vk, \overline{\word{i}}, \myvec{r}_i, \myvec{v}_i, \prfkzgwi)$. \newline
(3) Compute $c_i \leftarrow \rohash(\myvec{r}_i, \myvec{v}_i)$.
Retrieve $\barvec{r}$ and $\barvec{v}$ from $\pubstat$.
Assert $\veccomverify(\vk, \barvec{\myvec{r}} + c_i \cdot
	\barvec{\myvec{v}}, i, \myvec{r}_i + c_i \cdot \myvec{v}_i)$.



	


\caption{Main Framework for Generating Batch and Individual Proof}
\label{fig:main}
\end{algorithm*}
\end{scriptsize}
